% =============================================================================
% Handwritten Character Segmentation v2 --- Project Report
% Author:      Fahad Tariq (Matriculation: 23186464)
% Supervisor:  Dr.-Ing. Vincent Christlein
% Institution: Pattern Recognition Lab, FAU Erlangen-Nürnberg
% =============================================================================

\documentclass[a4paper,11pt]{article}

\usepackage[utf8]{inputenc}
\usepackage[T1]{fontenc}
\usepackage{lmodern}
\usepackage[british]{babel}
\usepackage{microtype}
\usepackage{amsmath,amssymb}
\usepackage{graphicx}
\usepackage{float}
\usepackage{caption}
\usepackage{subcaption}
\usepackage{booktabs}
\usepackage{multirow}
\usepackage{array}
\usepackage{tabularx}
\usepackage{longtable}
\usepackage[margin=2.5cm, headheight=14pt]{geometry}
\usepackage{fancyhdr}
\usepackage{enumitem}
\usepackage{xcolor}
\usepackage{xspace}
\usepackage{algorithm}
\usepackage{algpseudocode}
\usepackage[
    colorlinks=true,
    linkcolor=blue!65!black,
    citecolor=blue!65!black,
    urlcolor=blue!65!black,
    pdftitle={Handwritten Character Segmentation v2},
    pdfauthor={Fahad Tariq},
]{hyperref}

\pagestyle{fancy}
\fancyhf{}
\rhead{\small Handwritten Character Segmentation v2}
\lhead{\small Fahad Tariq --- FAU Erlangen-Nürnberg}
\cfoot{\thepage}
\renewcommand{\headrulewidth}{0.4pt}

\newcommand{\eg}{\emph{e.g.},\xspace}
\newcommand{\ie}{\emph{i.e.},\xspace}

% =============================================================================
\begin{document}
% =============================================================================

% -----------------------------------------------------------------------------
% Title Page
% -----------------------------------------------------------------------------
\begin{titlepage}
    \centering
    \vspace*{0.5cm}
    \includegraphics[width=0.55\textwidth]{fau_logo.png}

    \vfill

    {\LARGE\bfseries Handwritten Character Segmentation\par}
    \vspace{0.5cm}
    {\large Extended to Real Handwriting with Pixel-Level Masks\par}
    \vspace{1.5cm}
    {\large Project Report\par}

    \vfill

    \begin{tabular}{@{}ll@{}}
        Handed in by:         & Fahad Tariq                 \\[6pt]
        Matriculation number: & 23186464                    \\[6pt]
        Supervisor:           & Dr.-Ing. Vincent Christlein \\
    \end{tabular}

    \vspace*{\fill}
    {\small Pattern Recognition Lab
        \enspace---\enspace
        Martensstra\ss{}e~3
        \enspace---\enspace
        91058~Erlangen
        \enspace---\enspace
        lme.tf.fau.de\par}
    \vspace{0.5cm}
\end{titlepage}

% -----------------------------------------------------------------------------
% Abstract
% -----------------------------------------------------------------------------
\begin{abstract}
\noindent
This report presents the second iteration of a deep learning pipeline for handwritten
character segmentation, developed in direct response to supervisor feedback on the first
version. Three concrete improvements were made: segmentation masks are now generated at
the pixel level via a binarisation-and-AND procedure rather than filled bounding boxes,
ensuring that the model is trained to predict actual glyph contours rather than
rectangular regions; a render-time slant transform combined with an extended suite of
affine, elastic, morphological, and photometric augmentations is applied to better
simulate cursive and naturalistic handwriting variability; and trained models are
evaluated on real handwriting data from the IAM Handwriting Database for the first time,
providing a quantitative assessment of synthetic-to-real generalisation. Two architectures
are compared side by side---an Attention U-Net (29.3M parameters) and a SwinUNet based
on the Swin Transformer (41.4M parameters)---both trained on 15,000 synthetic images
spanning 80 character classes. On the synthetic test set, the Attention U-Net achieves a
best validation IoU of 98.08\% and a class-balanced macro IoU of 61.20\%, while SwinUNet
reaches 95.10\% validation IoU but a higher macro IoU of 68.20\%, indicating that
transformer-based global self-attention provides better per-class discrimination across a
long-tailed 80-class vocabulary. Evaluation on 1,861 real handwriting lines from the IAM
test split reveals a substantial and clearly two-tiered synthetic-to-real domain gap:
pixel-level accuracy degrades moderately to 84--87\%, confirming that low-level
ink-detection features partially transfer, whereas character-level macro IoU collapses to
below 0.3\% for both architectures, confirming that higher-level character-identity
features are tightly coupled to the font-rendered appearance of the synthetic domain and
do not generalise to natural handwriting. This dissociation precisely isolates the
remaining challenge for this task and motivates future work in domain adaptation, targeted
fine-tuning on real annotations, and stroke-level synthetic data generation.
\end{abstract}

\newpage
\tableofcontents
\newpage
\listoffigures
\listoftables
\newpage

% =============================================================================
\section{Introduction}
\label{sec:introduction}
% =============================================================================

\subsection{Background}

Segmenting individual characters in handwritten text is one of the harder sub-problems
in document image analysis. It sits directly upstream of recognition, so errors here
propagate through any downstream OCR or text extraction system. What makes it genuinely
difficult is that handwriting is inherently personal: stroke width, slant, character
spacing, and ligature formation vary not just across writers but within a single line from
the same writer. Unlike printed text, there are no clean white gaps between characters to
exploit, and the boundaries between adjacent letters are often ambiguous even to a human
reader at the pixel level.

Deep learning, and convolutional semantic segmentation in particular, has changed the
landscape considerably. Architectures like U-Net~\cite{ronneberger2015} and its
attention-gated variant~\cite{oktay2018} learn to produce dense per-pixel class
predictions directly from data, bypassing the need for hand-engineered feature pipelines.
More recently, vision transformers have demonstrated that global self-attention can
capture long-range dependencies missed by purely local convolutional
operations~\cite{liu2021,cao2022}, making them worth evaluating on a task like character
segmentation where context across a word or line is meaningful.

\subsection{Summary of Version 1}

The first version of this project established a working baseline. It rendered
handwritten-style fonts on blank canvases, generated segmentation masks via template
matching, and trained a single Attention U-Net that achieved 91.31\% pixel-weighted IoU
on a synthetic test set. During supervisor review, three shortcomings were identified:

\begin{enumerate}[leftmargin=*, label=(\arabic*)]
    \item The segmentation masks were filled rectangular bounding boxes rather than
          pixel-accurate character contours. This meant the model was trained to predict
          boxes, not glyph shapes---a fundamental mismatch between the training objective
          and the desired output.
    \item No slant or cursive augmentation was applied, limiting the variety of writing
          styles seen during training and reducing robustness to real handwriting slant.
    \item Evaluation never left the synthetic domain. There was no assessment on real
          handwriting, so it was entirely unknown whether the model generalised at all to
          natural writing.
\end{enumerate}

\subsection{Objectives of Version 2}

This report addresses all three points directly, following explicit guidance from
Dr.-Ing.\ Vincent Christlein:

\begin{enumerate}[leftmargin=*, label=(\arabic*)]
    \item \textbf{Pixel-level masks.} Replace bounding-box fills with glyph-contour masks
          generated by rendering each character in isolation, binarising the result, and
          taking the pixel-wise AND with the character region. The resulting masks follow
          actual ink shapes, including thin strokes, curves, and descenders.
    \item \textbf{Augmentation pipeline.} Introduce render-time horizontal shear to
          simulate cursive slant, and add a comprehensive training-time augmentation
          pipeline covering affine transforms, elastic deformation, morphological erosion
          and dilation (simulating pen pressure variation), photometric changes, and
          Gaussian noise (simulating scanning and paper texture artefacts).
    \item \textbf{Real data evaluation.} Evaluate trained models on the IAM Handwriting
          Database~\cite{marti2002} and produce a quantitative domain gap analysis that
          compares synthetic and real performance side by side.
\end{enumerate}

As an additional contribution beyond the three requested improvements, two architectures
are trained and compared---Attention U-Net and SwinUNet---to assess whether
transformer-based global context modelling offers advantages over the convolutional
approach for this task.

\subsection{Contributions}

The main contributions of this work are:

\begin{enumerate}[leftmargin=*, label=(\arabic*)]
    \item A pixel-accurate mask generation pipeline producing true glyph-contour
          annotations via binarisation and pixel-wise AND operations, correcting the
          fundamental training objective mismatch in version~1.
    \item A render-time slant augmentation scheme and an extended training-time pipeline
          combining geometric, morphological, photometric, and noise transforms, each
          motivated by a specific property of real handwriting variability.
    \item A controlled comparison between Attention U-Net and SwinUNet architectures on
          the same 80-class segmentation task, revealing a meaningful trade-off between
          pixel-level boundary precision and class-level discrimination.
    \item The first quantitative domain gap analysis for this system, evaluated on 1,861
          real handwriting images from the IAM test split, with a two-tiered
          characterisation separating low-level feature transfer from high-level class
          generalisation failure.
    \item A modular, reproducible codebase with full experiment tracking via
          Weights~\&~Biases~\cite{wandb2020}, centralised YAML configuration, and separate
          evaluation scripts for synthetic and real data.
\end{enumerate}

\subsection{Report Structure}

Section~\ref{sec:methodology} covers the full methodology: dataset generation, the pixel
mask procedure, the augmentation pipeline, both model architectures, training
configuration, and the IAM evaluation protocol. Section~\ref{sec:results} presents
quantitative and qualitative results on both synthetic and real data.
Section~\ref{sec:discussion} discusses what the results mean, where the current approach
falls short, and what future directions are most promising.
Section~\ref{sec:conclusion} summarises the work.

% =============================================================================
\section{Methodology}
\label{sec:methodology}
% =============================================================================

\subsection{Synthetic Dataset Generation}

\subsubsection{Text Sampling and Canvas Rendering}

Text content is sampled from Wikipedia articles using the \texttt{wikipedia-api} library,
which provides broad linguistic diversity across topics, vocabulary, and punctuation
patterns. Each image is produced by rendering 3--7 lines of text onto a $512 \times 512$
pixel white canvas (intensity 255). A font is selected at random for each image from a
curated collection of handwritten-style TrueType fonts sourced from Google Fonts, at a
randomly chosen size between 28 and 48 pixels. Text intensity is sampled uniformly
between 0 and 60 on a grey scale, introducing ink darkness variation that simulates
differences in pen ink density and paper contrast. The complete synthetic dataset contains
15,000 images partitioned into 12,000 training, 1,500 validation, and 1,500 test images.

\subsubsection{Pixel-Level Mask Generation}
\label{sec:pixel_masks}

The central correction from version~1 is the mask generation procedure. Previously, each
character was assigned a filled rectangular bounding box in the mask, which told the
model to predict boxes rather than character shapes. This is a fundamental training
objective mismatch: the model optimises its predictions to match rectangular fills, but
the desired output for a useful segmentation system is a mask that follows the actual ink
contour of each character.

The new approach generates a true glyph-contour mask for every character:

\begin{enumerate}[leftmargin=*, label=\arabic*.]
    \item Render the character in isolation on a blank white canvas at its computed layout
          position, using the same font, size, and colour as in the composite image.
    \item Binarise the rendered glyph with a fixed threshold of 128, producing a binary
          mask where ink pixels are 1 and background pixels are 0.
    \item Write the character's class index into the segmentation mask \emph{only at
          foreground (ink) pixel locations}---the pixel-wise AND of the bounding region
          with the binarised glyph.
\end{enumerate}

This procedure means thin strokes produce thin masks, curves produce curved masks, and
descenders extend below the baseline exactly as they appear in the rendered image. The
result is a ground-truth mask that is spatially co-registered with the visual ink, making
the segmentation task well-defined at the pixel level. When two characters overlap (rare
in practice with the chosen font sizes and spacing), the later character's class index
overwrites the earlier one at shared pixels, which mirrors the visual rendering order.

\begin{figure}[htbp]
    \centering
    \includegraphics[width=\textwidth]{figures/dataset_samples.png}
    \caption{Representative samples from the synthetic training dataset. Each row shows
    the rendered input image (left), the pixel-level ground-truth segmentation mask
    (centre), and a colour overlay blending both (right). Each colour in the mask
    corresponds to a distinct character class. The masks follow actual glyph contours
    rather than rectangular bounding boxes, meaning thin strokes, curves, descenders,
    and inter-character spacing are all labelled at the pixel level. Note the diversity
    of handwriting styles across rows, introduced by using multiple handwritten-style
    fonts at varying sizes and slant angles.}
    \label{fig:dataset_samples}
\end{figure}

\subsubsection{Character Vocabulary}

The vocabulary spans 80 classes in total: one background class (index 0) plus 79
character classes. These comprise 26 uppercase letters (A--Z), 26 lowercase letters
(a--z), 10 digits (0--9), and 17 special characters including space, period, comma,
semicolon, colon, exclamation mark, question mark, apostrophe, hyphen, quotation marks,
parentheses, ampersand, slash, at sign, hash, and plus sign. The mapping is defined in a
single \texttt{charset.py} module imported by all other pipeline components, ensuring
there is one authoritative source for class indices throughout the codebase.

\subsection{Augmentation Pipeline}

Augmentation operates at two distinct stages: render time (applied while generating the
synthetic dataset) and training time (applied on-the-fly in the data loader). Keeping
them separate makes the generated dataset fully reproducible while still providing
stochastic diversity during training. Critically, every augmentation in this pipeline is
motivated by a specific property of real handwriting that is absent from clean
font-rendered images.

\subsubsection{Render-Time Slant Augmentation}

To simulate the forward or backward lean characteristic of cursive handwriting, a
horizontal shear transform is applied at render time with probability 0.5. The shear
factor $s$ is sampled uniformly from $[-0.4,\, 0.4]$, covering a wide range of slant
angles from strongly backward-leaning to strongly forward-leaning text. The shear
matrix, centred vertically at height $h$, is:

\begin{equation}
    \begin{bmatrix} x' \\ y' \end{bmatrix}
    =
    \begin{bmatrix} 1 & s \\ 0 & 1 \end{bmatrix}
    \begin{bmatrix} x \\ y \end{bmatrix}
    +
    \begin{bmatrix} -sh/2 \\ 0 \end{bmatrix}
    \label{eq:shear}
\end{equation}

Critically, the \emph{same} shear matrix is applied to both the rendered image (using
bilinear interpolation) and the segmentation mask (using nearest-neighbour interpolation
to preserve integer class indices). This is essential for mask consistency: applying
shear to the image without applying the same transform to the mask would produce
misaligned supervision signal and corrupt the training objective.

\subsubsection{Training-Time Augmentations}

Additional augmentations are applied stochastically during data loading via the
Albumentations library~\cite{albumentations2020}, combined with two custom transforms
for morphological operations. The pipeline is organised into four groups. The rationale
for each group is stated alongside the parameters, connecting each augmentation to the
specific real handwriting property it targets.

\medskip
\noindent\textbf{Geometric transforms} (applied jointly to image and mask, probability
0.5):
\begin{itemize}[noitemsep]
    \item Affine: rotation ($\pm 3^{\circ}$), shear ($\pm 5^{\circ}$), scale (0.95--1.05), translation
          ($\pm 3\%$). \emph{Rationale:} real writers do not place text on a perfectly
          horizontal baseline; small global misalignments are ubiquitous in scanned
          documents.
    \item Perspective distortion (scale 0.02--0.05). \emph{Rationale:} documents are
          rarely scanned flat; perspective foreshortening introduces systematic spatial
          distortion.
    \item Elastic deformation ($\alpha = 30$, $\sigma = 5$). \emph{Rationale:} the
          waviness and local curvature of hand-drawn strokes arise from natural motor
          tremor and pen velocity variation that font rendering cannot replicate.
\end{itemize}

\medskip
\noindent\textbf{Morphological transforms} (image only, probability 0.3):
\begin{itemize}[noitemsep]
    \item Erosion (kernel up to $2 \times 2$): thins ink strokes. \emph{Rationale:} light
          pen pressure or fast writing produces thin, faint strokes that erode the
          character body.
    \item Dilation (kernel up to $2 \times 2$): thickens ink strokes. \emph{Rationale:}
          heavy pen pressure or slow writing produces broad, dense strokes that expand the
          character body.
    \item Simulated ink variation: localised brightness perturbation on foreground pixels
          only. \emph{Rationale:} ink flow in a real pen is not constant; starved or
          flooded regions produce within-stroke intensity gradients absent from rendered
          text.
\end{itemize}

\noindent Morphological transforms are applied only to the image, not the mask, because
they alter the visual appearance of ink strokes without changing the semantic identity of
the character being written.

\medskip
\noindent\textbf{Photometric transforms} (image only, probability 0.5):
\begin{itemize}[noitemsep]
    \item Random brightness and contrast ($\pm 15\%$). \emph{Rationale:} paper yellowing,
          scanner calibration differences, and ambient lighting all affect document image
          contrast.
    \item CLAHE (clip limit 2.0). \emph{Rationale:} local histogram equalisation simulates
          the contrast normalisation often applied during document pre-processing, exposing
          the model to inputs that may arrive pre-processed at inference time.
\end{itemize}

\medskip
\noindent\textbf{Noise transforms} (image only, probability 0.3):
\begin{itemize}[noitemsep]
    \item Gaussian noise (variance up to 15.0). \emph{Rationale:} CCD sensor noise in
          scanner hardware introduces additive Gaussian noise, particularly in lightly
          illuminated regions.
    \item Gaussian blur (kernel size up to $3 \times 3$). \emph{Rationale:} out-of-focus
          scanning and ink bleed into paper fibres both produce low-pass blurring of stroke
          edges.
\end{itemize}

\medskip
\noindent Geometric transforms must be applied to both image and mask simultaneously to
preserve label alignment. All other transforms affect only pixel intensities and are
therefore applied to the image alone.

\subsection{Model Architectures}

\subsubsection{Attention U-Net}

The Attention U-Net~\cite{oktay2018} follows a standard encoder-decoder structure with
skip connections between corresponding encoder and decoder levels, enhanced with attention
gates. The encoder consists of four blocks with channel dimensions
$[64, 128, 256, 512]$, each containing two $3 \times 3$ convolutions with batch
normalisation and ReLU activation, followed by $2 \times 2$ max pooling for
downsampling. A bottleneck block operates at 1,024 channels.

The decoder comprises four symmetric upsampling blocks. Bilinear upsampling followed by
a $1 \times 1$ projection is used in preference to transposed convolution, which is prone
to checkerboard artefacts. Each decoder block applies an attention gate to the incoming
skip connection before concatenation. The gate computes attention coefficients from both
the decoder (gating) signal $\mathbf{g}$ and the encoder skip signal $\mathbf{x}^l$:

\begin{equation}
    \alpha^l = \sigma_2\!\left( \mathbf{W}_\psi \cdot \sigma_1\!\left(
    \mathbf{W}_x \mathbf{x}^l + \mathbf{W}_g \mathbf{g} + b_g \right)
    + b_\psi \right)
    \label{eq:attention_gate}
\end{equation}

where $\sigma_1$ is ReLU, $\sigma_2$ is sigmoid, and $\mathbf{W}_x$, $\mathbf{W}_g$,
$\mathbf{W}_\psi$ are learned linear projections. The resulting attention map
$\alpha^l \in [0,1]$ is applied element-wise to $\mathbf{x}^l$, suppressing background
regions and amplifying character boundary features before the skip connection reaches the
decoder. Dropout ($p = 0.1$) is applied after each decoder block. The output is a
$1 \times 1$ convolution producing 80 channel logits. Total trainable parameters:
\textbf{29.3M}.

\subsubsection{SwinUNet}

The SwinUNet~\cite{cao2022} replaces convolutional blocks throughout the encoder and
decoder with Swin Transformer blocks~\cite{liu2021}, making it a pure-transformer
segmentation architecture. Input images are first partitioned into non-overlapping
$4 \times 4$ patches and projected to an embedding dimension of 96. Four encoder stages
process these tokens with depths $[2, 2, 6, 2]$ and attention head counts
$[3, 6, 12, 24]$.

Within each stage, Swin Transformer blocks alternate between standard window attention
and shifted-window attention. Regular window attention restricts self-attention to
non-overlapping $8 \times 8$ windows, reducing quadratic complexity from
$\mathcal{O}(n^2)$ to $\mathcal{O}(n \cdot w^2)$ where $w$ is the window size.
Shifted-window attention in alternating blocks displaces window boundaries by half a
window size, enabling cross-window information exchange without global attention.
Learnable relative position biases provide spatial awareness within each window.

Patch merging between encoder stages halves the spatial resolution and doubles the
channel dimension, analogous to max pooling in a CNN. The decoder mirrors this with patch
expanding layers for upsampling and linear projections for skip connection fusion.
Stochastic depth regularisation applies a linearly increasing drop probability from 0 to
0.1 across layers. The final projection maps to 80 output classes. Total trainable
parameters: \textbf{41.4M}.

The key architectural distinction relevant to this task is that global self-attention
across the full sequence of patch tokens is gradually accumulated through the
shifted-window mechanism, allowing the model to relate spatially distant characters within
a line. This is qualitatively different from the local receptive field growth of a CNN
encoder, and may contribute to better discrimination between visually similar character
classes when their surrounding context is informative.

\subsection{Training Procedure}

\subsubsection{Loss Function}

Training minimises a combined loss~\cite{sudre2017}:
\begin{equation}
    \mathcal{L}_{\text{total}} = 0.5 \cdot \mathcal{L}_{\text{Dice}}
    + 0.5 \cdot \mathcal{L}_{\text{CE}}
    \label{eq:loss}
\end{equation}

The Dice component is computed per class and averaged \emph{only over classes present in
the current batch}, avoiding the degenerate case where absent classes contribute a
near-zero-prediction Dice score of approximately 1.0 that would dominate the gradient
with an uninformative signal:

\begin{equation}
    \mathcal{L}_{\text{Dice}} = 1 - \frac{1}{|C_{\text{present}}|}
    \sum_{c \in C_{\text{present}}}
    \frac{2\displaystyle\sum_i p_{i,c}\, g_{i,c} + \epsilon}
         {\displaystyle\sum_i p_{i,c} + \sum_i g_{i,c} + \epsilon}
    \label{eq:dice}
\end{equation}

where $C_{\text{present}}$ is the set of classes with at least one ground-truth pixel in
the batch, $p_{i,c}$ and $g_{i,c}$ are the predicted probability and one-hot ground
truth for pixel $i$ and class $c$, and $\epsilon = 1.0$ is a smoothing constant.

The Cross-Entropy component uses per-class weights to address the large class-frequency
imbalance across 80 classes. Weights are computed as inverse class frequency, then
dampened by a square-root transform to avoid extreme values for very rare classes, and
clamped to a maximum of 10.0. The background class weight is further reduced by 95\% so
that the abundant background pixels do not dominate the gradient. Label smoothing of 0.01
prevents overconfident predictions.

The combination is motivated by complementary strengths: Cross-Entropy provides stable,
well-behaved gradients throughout training and handles multi-class classification
naturally; Dice loss directly optimises the overlap metric that is reported at evaluation
time and is inherently robust to class imbalance by normalising over predicted and
ground-truth areas.

\subsubsection{Optimisation and Training Configuration}

All models are trained with the AdamW optimiser~\cite{loshchilov2019} at a learning rate
of $3 \times 10^{-4}$ and weight decay $1 \times 10^{-4}$. Cosine annealing reduces the
learning rate to a minimum of $1 \times 10^{-6}$ over 100 epochs. Gradient norms are
clipped to 1.0. Training uses 16-bit mixed precision and PyTorch DDP across 4 NVIDIA GPUs
with per-GPU batch size 4 and gradient accumulation over 4 steps, giving an effective
batch size of 64. Early stopping monitors validation IoU with patience of 25 epochs.
Training is managed by PyTorch Lightning~\cite{falcon2019}. Table~\ref{tab:training_config}
summarises the key hyperparameters.

\begin{table}[H]
\centering
\caption{Training Configuration}
\label{tab:training_config}
\begin{tabular}{ll}
\toprule
\textbf{Hyperparameter} & \textbf{Value} \\
\midrule
Optimiser            & AdamW \\
Learning rate        & $3 \times 10^{-4}$ \\
Weight decay         & $1 \times 10^{-4}$ \\
LR schedule          & Cosine annealing ($\eta_{\min} = 10^{-6}$, $T_{\max} = 100$) \\
Effective batch size & 64 (4/GPU $\times$ 4 GPUs $\times$ 4 accum.) \\
Loss                 & Dice + CE (0.5 each), present-class averaging \\
Label smoothing      & 0.01 \\
Gradient clipping    & max norm 1.0 \\
Precision            & 16-bit mixed \\
Early stopping       & patience 25, monitor val/IoU (macro, pixel-weighted) \\
\bottomrule
\end{tabular}
\end{table}

\subsection{IAM Evaluation Protocol}

\subsubsection{Dataset}

The IAM Handwriting Database~\cite{marti2002} contains English text handwritten by 657
writers on pre-printed forms. We evaluate on 1,861 line images from the standard Large
Writer Independent Text Line Recognition Task test split. IAM provides word-level
bounding box annotations and character-level transcriptions, but does not provide
pixel-level segmentation masks. To enable quantitative evaluation without manual
annotation, pseudo ground-truth masks are generated automatically.

\subsubsection{Pseudo Mask Generation}
\label{sec:pseudo_masks}

Since pixel-level character segmentation annotations are unavailable in IAM, approximate
masks are constructed as follows:

\begin{enumerate}[leftmargin=*, label=\arabic*.]
    \item Apply Otsu's automatic thresholding to binarise the handwriting image, adapting
          to each image's local contrast.
    \item Extract connected components with 8-connectivity and discard any with fewer than
          10 pixels (noise suppression).
    \item Group components into text lines by vertical centroid proximity.
    \item Within each line, sort components left-to-right by horizontal centroid.
    \item Assign components to transcription characters sequentially, mapping each
          component's pixels to the class index of the corresponding character.
\end{enumerate}

\noindent\textbf{Edge case handling.} When the number of extracted components does not
match the number of transcription characters, the following rules apply: if there are
\emph{fewer} components than characters (most commonly due to touching characters merging
into a single connected region), the available components are assigned to the first $k$
transcription characters and remaining characters receive no mask pixels. If there are
\emph{more} components than characters (most commonly due to dotted or multi-stroke
letters such as `i', `j', or `t' producing multiple disconnected regions), excess
components beyond the transcription length are discarded and assigned to background.

This procedure is necessarily approximate. Touching characters, split characters, and
reading-order ambiguities all introduce systematic noise into the resulting masks.
Real cursive handwriting is particularly problematic: characters frequently connect
into a single continuous pen stroke, causing multiple characters to merge into one
connected component that is then incorrectly attributed to only the first transcription
character in the group. Consequently, IAM metrics \emph{underestimate} true model
performance to some degree and should be interpreted as a relative lower bound for domain
gap assessment rather than an absolute accuracy figure. The primary purpose of the IAM
evaluation is to establish that a domain gap exists and to quantify its structural
character, not to report production-quality accuracy on real handwriting.

\subsection{Evaluation Metrics}
\label{sec:metrics}

All metrics exclude the background class (index 0) from macro averaging. Three levels of
granularity are reported:

\medskip
\noindent\textbf{Overall metrics} (macro-averaged over active foreground classes):
\begin{align}
    \text{IoU}_c &= \frac{|\hat{P}_c \cap G_c|}{|\hat{P}_c \cup G_c|}, \qquad
    \overline{\text{IoU}} = \frac{1}{|C_{\text{active}}|}
    \sum_{c \in C_{\text{active}}} \text{IoU}_c
    \label{eq:iou} \\[6pt]
    \text{Dice}_c &= \frac{2|\hat{P}_c \cap G_c|}{|\hat{P}_c| + |G_c|}, \qquad
    \text{Pixel Acc} = \frac{\sum_c \text{TP}_c}{\text{total pixels}}
    \label{eq:dice_eval}
\end{align}

\noindent\textbf{Per-group metrics}: Aggregated separately for uppercase, lowercase,
digits, and special characters.

\noindent\textbf{Per-class metrics}: Individual IoU, Dice, precision, recall, and F1 for
each of the 79 foreground classes, enabling fine-grained analysis of which characters are
hardest to segment.

\medskip
\noindent\textbf{Important note on metric computation.} Two distinct IoU computations
appear in this report and must not be conflated. During training, validation IoU is
computed online by \texttt{torchmetrics} \texttt{MulticlassJaccardIndex} with
\texttt{average="macro"}, which averages per-class IoU values weighted by the number of
pixels per class---this gives larger, more frequent classes greater influence and drives
the early stopping criterion. At evaluation time, the full confusion matrix is accumulated
over the entire test set and macro IoU is computed by giving every foreground class
\emph{equal weight} regardless of pixel frequency---a class-balanced average that
penalises underperformance on rare characters equally with common ones.

These two quantities are not directly comparable, and a model may rank differently on
each. In particular, a model can achieve high pixel-weighted validation IoU by excelling
on high-frequency classes while failing on rare ones, yet score substantially lower on the
class-balanced test macro IoU. The class-balanced macro IoU is the primary evaluation
metric reported in all result tables throughout this report, as it provides a more
demanding and informative measure of segmentation quality across the full 80-class
vocabulary.

A direct consequence for model selection is that early stopping based on pixel-weighted
validation IoU does not guarantee selection of the checkpoint that maximises
class-balanced test IoU. This metric inconsistency is discussed further in
Section~\ref{sec:metric_inconsistency}.

\subsection{Implementation Details}

The system is implemented in Python 3.11 with PyTorch 2.1~\cite{paszke2019} (CUDA 11.8).
PyTorch Lightning~2.1.4~\cite{falcon2019} manages the training loop, distributed
training, and checkpointing. Experiment tracking and hyperparameter visualisation use
Weights~\&~Biases~\cite{wandb2020}. Augmentations are implemented with
Albumentations~1.4.4~\cite{albumentations2020}, extended with two custom
\texttt{ImageOnlyTransform} classes for morphological erosion and dilation. Configuration
is centralised in a single YAML file parsed by OmegaConf, making every experimental
setting reproducible from one file. All experiments use random seed 42 for
reproducibility.

% =============================================================================
\section{Results}
\label{sec:results}
% =============================================================================

\subsection{Training Dynamics}

\subsubsection{Attention U-Net}

Training ran for 82 epochs before early stopping triggered, with the best checkpoint
identified at epoch~67 (validation IoU 98.08\%, extracted from the saved
\texttt{.ckpt} file). Learning was rapid in the early phase: validation IoU climbed from
18.4\% at epoch~0 to 64.0\% by epoch~6 and crossed 93\% by epoch~28, indicating that
the attention mechanism and pixel-level masks together provided a clean, strong training
signal from the outset. Each epoch completed in approximately 3~minutes 14~seconds
across the four GPUs.

\subsubsection{SwinUNet}

SwinUNet converged more slowly in the first few epochs---starting from 0.5\% validation
IoU at epoch~0 compared to 18.4\% for the Attention U-Net. This slow start is consistent
with transformer architectures requiring more gradient updates to initialise meaningful
attention patterns from random weights. After this warm-up phase, SwinUNet converged
steadily, with the best checkpoint identified at epoch~61 (validation IoU 95.10\%,
extracted from the saved \texttt{.ckpt} file). Early stopping halted training at
epoch~86. Despite having 41\% more parameters, each epoch completed in approximately
3~minutes, marginally faster than the Attention U-Net, reflecting the more parallelisable
nature of transformer matrix operations on GPU hardware.

\subsubsection{Training Summary}

Table~\ref{tab:training_summary} provides a compact summary of the training run for both
architectures for direct comparison.

\begin{table}[H]
\centering
\caption{Training Run Summary}
\label{tab:training_summary}
\begin{tabular}{lcccccc}
\toprule
\textbf{Architecture} & \textbf{Best Epoch} & \textbf{Best Val IoU} & \textbf{Best Val Dice}
    & \textbf{Total Epochs} & \textbf{Time / Epoch} \\
\midrule
Attention U-Net & 67 & 98.08\% & 96.10\% & 82 & $\approx$3 min 14 s \\
SwinUNet        & 61 & 95.10\% & 92.81\% & 86 & $\approx$3 min 00 s \\
\bottomrule
\multicolumn{6}{l}{\footnotesize Best-epoch values are extracted from saved \texttt{.ckpt}
files; see Figure~\ref{fig:training_curves} for details.} \\
\end{tabular}
\end{table}

\begin{figure}[htbp]
    \centering
    \begin{subfigure}[t]{\textwidth}
        \centering
        \includegraphics[width=\textwidth]{figures/unet_training_curves.png}
        \caption{Attention U-Net: validation IoU (left) and validation Dice (right) over
        training epochs.}
        \label{fig:unet_training_curves}
    \end{subfigure}

    \vspace{1em}

    \begin{subfigure}[t]{\textwidth}
        \centering
        \includegraphics[width=\textwidth]{figures/swin_training_curves.png}
        \caption{SwinUNet: validation IoU (left) and validation Dice (right) over training
        epochs.}
        \label{fig:swin_training_curves}
    \end{subfigure}

    \caption{Training dynamics for both architectures. Validation IoU and validation Dice
    are plotted over the epochs recorded by the W\&B logger. The grey shaded region in
    each plot (approximately epochs 0--45) corresponds to early training during which the
    W\&B run was interrupted by an HPC job preemption; these epochs completed successfully
    and their model checkpoints were saved locally, but metrics were not synced to the
    W\&B server. The \textbf{Best} annotations visible within each subplot indicate the
    peak values observed within the logged portion of training only (post-epoch~45)---they
    are \emph{not} the global best values across the full run. The authoritative
    best-checkpoint values---\textbf{98.08\%} validation IoU at epoch~67 for the Attention
    U-Net and \textbf{95.10\%} at epoch~61 for SwinUNet---were identified from the saved
    \texttt{.ckpt} files by reading the \texttt{best\_model\_score} field recorded by the
    \texttt{ModelCheckpoint} callback at the time of saving, independently of W\&B. All
    result tables in this report use these checkpoint-derived values; see
    Table~\ref{tab:training_summary}.}
    \label{fig:training_curves}
\end{figure}

\subsection{Synthetic Test Set Evaluation}

\subsubsection{Overall Metrics}

Table~\ref{tab:synthetic_results} reports the full evaluation on the held-out synthetic
test set of 1,500 images. Bold entries indicate the better score for each metric.

\begin{table}[H]
\centering
\caption{Synthetic Test Set Performance Comparison}
\label{tab:synthetic_results}
\begin{tabular}{lcc}
\toprule
\textbf{Metric} & \textbf{Attention U-Net} & \textbf{SwinUNet} \\
\midrule
Best Validation IoU (pixel-weighted)$^\dagger$ & \textbf{98.08\%} & 95.10\% \\
Pixel Accuracy                                 & 98.59\%          & \textbf{99.03\%} \\
Macro IoU (class-balanced)                     & 61.20\%          & \textbf{68.20\%} \\
Macro Dice (class-balanced)                    & 75.07\%          & \textbf{80.32\%} \\
Precision                                      & 61.30\%          & \textbf{68.61\%} \\
Recall                                         & \textbf{99.63\%} & 98.91\% \\
F1 Score                                       & 75.90\%          & \textbf{81.02\%} \\
Parameters                                     & \textbf{29.3M}   & 41.4M \\
\bottomrule
\multicolumn{3}{l}{$^\dagger$ Computed by torchmetrics during training (pixel-weighted);
see Section~\ref{sec:metrics}.} \\
\end{tabular}
\end{table}

The results reveal an important split between the two models. The Attention U-Net achieves
higher validation IoU (98.08\% versus 95.10\%), yet SwinUNet outperforms it on every
macro-averaged test metric. This discrepancy arises from the difference between the two
IoU computations described in Section~\ref{sec:metrics}. Validation IoU is
pixel-weighted, so larger classes contribute more to the score; this benefits U-Net's
strong local boundary precision on high-frequency characters. The class-balanced macro IoU
computed at test time gives equal weight to every foreground class regardless of
frequency, which penalises failure on rare characters equally with common ones.

The fact that SwinUNet wins on all class-balanced test metrics despite losing on
pixel-weighted validation IoU suggests that its global self-attention mechanism helps it
discriminate between visually confusable characters---such as `O' vs.\ `0', `l' vs.\ `1',
and `I' vs.\ `|'---that appear rarely and are essentially invisible in a pixel-weighted
average but heavily penalised in a class-balanced one. The F1 scores (75.90\% for U-Net,
81.02\% for SwinUNet) are consistent with this interpretation: the harmonic mean of
precision and recall is sensitive to failures on both under-predicted and over-predicted
classes, and SwinUNet's advantage of 5.12 percentage points on F1 mirrors its advantage
on macro IoU and Dice, confirming a systematic improvement in class-level discrimination
rather than a single-metric artefact. This finding has a methodological implication: for
a long-tailed 80-class vocabulary, class-balanced macro IoU is a more informative
evaluation metric than pixel-weighted IoU, and optimising directly for it (rather than
using it only at test time) could further improve SwinUNet's advantage.

Table~\ref{tab:v1v2} places version~2 in context by comparing the Attention U-Net against
the version~1 baseline on aspects that can be meaningfully compared. The test-set IoU
figures from the two versions are presented separately because they measure fundamentally
different quantities and cannot be placed in the same column without misleading the
reader: version~1 reported pixel-weighted IoU over bounding-box ground truth, while
version~2 reports class-balanced macro IoU over glyph-contour ground truth with 80
classes. The validation IoU figures are the only directly comparable numbers, as both use
the same torchmetrics pixel-weighted computation on their respective validation sets.

\begin{table}[H]
\centering
\caption{Version 1 vs.\ Version 2: Attention U-Net --- Pipeline Comparison}
\label{tab:v1v2}
\begin{tabular}{lcc}
\toprule
\textbf{Aspect} & \textbf{Version 1} & \textbf{Version 2} \\
\midrule
Mask type              & Bounding-box fill       & Glyph-contour (binarise + AND) \\
Character classes      & 75                      & 80 \\
Training images        & 8,000                   & 12,000 \\
Slant augmentation     & None                    & Render-time shear $[-0.4, 0.4]$ \\
Morphological augm.    & None                    & Erosion, dilation, ink variation \\
\midrule
\multicolumn{3}{l}{\textit{Directly comparable}} \\
\addlinespace[2pt]
Best validation IoU (pixel-weighted) & 94.52\% & \textbf{98.08\%} \\
\midrule
\multicolumn{3}{l}{\textit{Not directly comparable --- different ground truth and metric}} \\
\addlinespace[2pt]
Test IoU (v1: pixel-weighted, bbox GT) & 91.31\% & --- \\
Test IoU (v2: class-balanced, glyph GT) & --- & 61.20\% \\
\midrule
Real data evaluation   & None                    & IAM test split (1,861 images) \\
Architecture(s)        & Attention U-Net         & Attention U-Net + SwinUNet \\
\bottomrule
\end{tabular}
\end{table}

\subsubsection{Per-Group Analysis}

Table~\ref{tab:per_group} breaks the class-balanced macro IoU down by character group.
SwinUNet outperforms the Attention U-Net across all four groups without exception.

\begin{table}[H]
\centering
\caption{Per-Group IoU on Synthetic Test Set (Class-Balanced Macro)}
\label{tab:per_group}
\begin{tabular}{lcc}
\toprule
\textbf{Group} & \textbf{Attention U-Net} & \textbf{SwinUNet} \\
\midrule
Lowercase   & 76.38\% & \textbf{82.28\%} \\
Uppercase   & 54.94\% & \textbf{64.06\%} \\
Digits      & 49.51\% & \textbf{56.69\%} \\
Special     & 47.74\% & \textbf{51.78\%} \\
\bottomrule
\end{tabular}
\end{table}

Lowercase letters consistently achieve the highest IoU for both models, which follows
directly from their high frequency in English text and therefore in the training data:
more training examples per class leads to stronger per-class feature learning. Special
characters and digits are the most challenging: special characters include rare punctuation
with very few training samples, while digits suffer from structural similarity to letters
(\eg `1'/`l', `0'/`O'). The performance ordering (lowercase $>$ uppercase $>$ digits $>$
special) is consistent across both architectures and all three metrics, suggesting it
reflects a genuine property of the task rather than a random artefact.

\begin{figure}[htbp]
    \centering
    \begin{subfigure}[t]{0.48\textwidth}
        \centering
        \includegraphics[width=\textwidth]{figures/unet_synthetic_per_group_iou.png}
        \caption{U-Net --- per-group IoU}
        \label{fig:pg_syn_unet_iou}
    \end{subfigure}
    \hfill
    \begin{subfigure}[t]{0.48\textwidth}
        \centering
        \includegraphics[width=\textwidth]{figures/swin_synthetic_per_group_iou.png}
        \caption{SwinUNet --- per-group IoU}
        \label{fig:pg_syn_swin_iou}
    \end{subfigure}

    \vspace{0.6em}

    \begin{subfigure}[t]{0.48\textwidth}
        \centering
        \includegraphics[width=\textwidth]{figures/unet_synthetic_per_group_dice.png}
        \caption{U-Net --- per-group Dice}
        \label{fig:pg_syn_unet_dice}
    \end{subfigure}
    \hfill
    \begin{subfigure}[t]{0.48\textwidth}
        \centering
        \includegraphics[width=\textwidth]{figures/swin_synthetic_per_group_dice.png}
        \caption{SwinUNet --- per-group Dice}
        \label{fig:pg_syn_swin_dice}
    \end{subfigure}

    \vspace{0.6em}

    \begin{subfigure}[t]{0.48\textwidth}
        \centering
        \includegraphics[width=\textwidth]{figures/unet_synthetic_per_group_f1.png}
        \caption{U-Net --- per-group F1}
        \label{fig:pg_syn_unet_f1}
    \end{subfigure}
    \hfill
    \begin{subfigure}[t]{0.48\textwidth}
        \centering
        \includegraphics[width=\textwidth]{figures/swin_synthetic_per_group_f1.png}
        \caption{SwinUNet --- per-group F1}
        \label{fig:pg_syn_swin_f1}
    \end{subfigure}

    \caption{Per-group segmentation metrics on the synthetic test set for both
    architectures across IoU (top), Dice (middle), and F1 (bottom). SwinUNet outperforms
    Attention U-Net across all groups and all three metrics, with the largest absolute gap
    in the uppercase group. The ordering lowercase $>$ uppercase $>$ digits $>$ special is
    consistent across all metrics and both models. All metrics are class-balanced macro
    averages within each group.}
    \label{fig:per_group_synthetic}
\end{figure}

\subsubsection{Qualitative Results on Synthetic Data}

Figure~\ref{fig:synthetic_predictions} shows representative model predictions on the
synthetic test set alongside ground-truth masks and colour overlays. Both architectures
produce visually accurate, well-aligned segmentations across diverse font styles and
character densities. Figure~\ref{fig:confusion_synthetic} presents the corresponding
normalised confusion matrices restricted to the top 30 most frequent classes. The strong
diagonal confirms overall high accuracy. The most notable off-diagonal entries are between
structurally similar pairs: `O'/`0', `l'/`1', and `I'/`|', which share visual features
that are difficult to disambiguate without strong contextual cues. SwinUNet's confusion
matrix shows a slightly stronger diagonal for these pairs, consistent with its better
macro IoU on digits and special characters.

\begin{figure}[htbp]
    \centering
    \begin{subfigure}[t]{\textwidth}
        \centering
        \includegraphics[width=\textwidth]{figures/unet_synthetic_predictions.png}
        \caption{Attention U-Net predictions on synthetic test samples.}
        \label{fig:unet_syn_pred}
    \end{subfigure}
    \vspace{0.6em}
    \begin{subfigure}[t]{\textwidth}
        \centering
        \includegraphics[width=\textwidth]{figures/swin_synthetic_predictions.png}
        \caption{SwinUNet predictions on synthetic test samples.}
        \label{fig:swin_syn_pred}
    \end{subfigure}
    \caption{Qualitative segmentation results on the synthetic test set. Each row shows
    the input image, ground-truth mask, model prediction, and a colour overlay. Both
    models achieve visually accurate character-level segmentation with precise
    glyph-contour boundaries on the synthetic domain, consistent with the high
    quantitative metrics in Table~\ref{tab:synthetic_results}.}
    \label{fig:synthetic_predictions}
\end{figure}

\begin{figure}[htbp]
    \centering
    \begin{subfigure}[t]{0.48\textwidth}
        \centering
        \includegraphics[width=\textwidth]{figures/unet_synthetic_confusion_matrix.png}
        \caption{Attention U-Net}
        \label{fig:cm_unet_syn}
    \end{subfigure}
    \hfill
    \begin{subfigure}[t]{0.48\textwidth}
        \centering
        \includegraphics[width=\textwidth]{figures/swin_synthetic_confusion_matrix.png}
        \caption{SwinUNet}
        \label{fig:cm_swin_syn}
    \end{subfigure}
    \caption{Normalised confusion matrices on the synthetic test set (top 30 classes by
    frequency). Rows are true classes, columns are predicted classes. The strong diagonal
    confirms high overall accuracy. Systematic off-diagonal confusion is concentrated on
    visually similar pairs: `O'/`0', `l'/`1', and `I'/`|'. SwinUNet shows marginally
    stronger diagonal entries for these ambiguous pairs, consistent with its higher macro
    IoU on digits and special characters.}
    \label{fig:confusion_synthetic}
\end{figure}

\subsection{IAM Handwriting Database Evaluation}

\subsubsection{Quantitative Results}
\label{sec:iam_quantitative}

Table~\ref{tab:iam_results} presents the evaluation results on the IAM test split. The
sharp contrast with Table~\ref{tab:synthetic_results} is the defining result of this
work.

\begin{table}[H]
\centering
\caption{IAM Test Set Performance (1,861 Line Images, Pseudo Ground Truth)}
\label{tab:iam_results}
\begin{tabular}{lcc}
\toprule
\textbf{Metric} & \textbf{Attention U-Net} & \textbf{SwinUNet} \\
\midrule
Pixel Accuracy & \textbf{87.01\%} & 84.48\% \\
Macro IoU      & 0.16\%           & \textbf{0.26\%} \\
Macro Dice     & 0.32\%           & \textbf{0.52\%} \\
Precision      & 0.85\%           & \textbf{0.86\%} \\
Recall         & \textbf{1.30\%}  & 1.01\% \\
\bottomrule
\multicolumn{3}{l}{\footnotesize Metrics use approximate pseudo ground-truth masks; see
Section~\ref{sec:pseudo_masks}. Values underestimate true performance.} \\
\end{tabular}
\end{table}

Pixel accuracy holds up reasonably well at 84--87\%, confirming that both models retain
the ability to distinguish ink pixels from background on real handwriting images.
Character-level discrimination, however, collapses entirely: macro IoU falls to
0.16--0.26\%. To contextualise this figure, a uniform random classifier assigning one of
79 foreground classes to each ink pixel would achieve an expected macro IoU of
approximately $1/79 \approx 1.27\%$. Both models score \emph{below} this random baseline,
which is not a sign of random failure but rather a sign of systematic class-bias: the
models predict predominantly high-frequency lowercase classes for virtually all ink pixels,
yielding near-zero IoU for the many classes they essentially never predict. This behaviour
is directly visible in Figure~\ref{fig:iam_predictions}: the overlay column confirms that
foreground ink regions are detected in roughly the correct spatial locations, but the
colour-coded class labels assigned to those ink pixels do not match the true character
identity. The models have learned where ink is, but not which character that ink
represents when the handwriting domain departs from the synthetic training distribution.

\subsubsection{Domain Gap Analysis}

Table~\ref{tab:domain_gap} quantifies the performance degradation from the synthetic to
the real domain for both architectures.

\begin{table}[H]
\centering
\caption{Domain Gap: Synthetic Test Set vs.\ IAM Performance}
\label{tab:domain_gap}
\begin{tabular}{lccc}
\toprule
\textbf{Metric} & \textbf{Synthetic} & \textbf{IAM} & \textbf{Relative Drop} \\
\midrule
\multicolumn{4}{l}{\textit{Attention U-Net}} \\
\addlinespace[2pt]
Pixel Accuracy & 98.59\% & 87.01\% & $-11.7\%$ \\
Macro IoU      & 61.20\% & 0.16\%  & $-99.7\%$ \\
Macro Dice     & 75.07\% & 0.32\%  & $-99.6\%$ \\
\midrule
\multicolumn{4}{l}{\textit{SwinUNet}} \\
\addlinespace[2pt]
Pixel Accuracy & 99.03\% & 84.48\% & $-14.7\%$ \\
Macro IoU      & 68.20\% & 0.26\%  & $-99.6\%$ \\
Macro Dice     & 80.32\% & 0.52\%  & $-99.3\%$ \\
\bottomrule
\end{tabular}
\end{table}

The two-tier degradation pattern is the central empirical finding of this work. Pixel
accuracy drops by 12--15\%---substantial but recoverable---while class-level metrics drop
by over 99\% in both cases. This dissociation directly reveals the structure of what is
and is not learned from synthetic font-rendered training data: low-level ink-detection
features (intensity gradients, stroke edges) are general enough to transfer to real
handwriting, but the higher-level character-identity features are so deeply coupled to
the smooth, consistent appearance of font-rendered glyphs that they cannot recognise the
same characters when drawn with natural pen variation.

\begin{figure}[htbp]
    \centering
    \begin{subfigure}[t]{0.48\textwidth}
        \centering
        \includegraphics[width=\textwidth]{figures/unet_domain_gap_chart.png}
        \caption{Attention U-Net}
        \label{fig:gap_unet}
    \end{subfigure}
    \hfill
    \begin{subfigure}[t]{0.48\textwidth}
        \centering
        \includegraphics[width=\textwidth]{figures/swin_domain_gap_chart.png}
        \caption{SwinUNet}
        \label{fig:gap_swin}
    \end{subfigure}
    \caption{Domain gap visualisation: synthetic test set vs.\ IAM performance across all
    metrics for both architectures. The two-tier structure is clearly visible: pixel
    accuracy degrades moderately while all character-class metrics collapse to near zero,
    confirming that the domain gap affects class-identity features far more severely than
    low-level ink-detection features.}
    \label{fig:domain_gap_charts}
\end{figure}

\subsubsection{Per-Group Domain Gap}

Tables~\ref{tab:group_gap} and~\ref{tab:group_gap_swin} show the per-group IoU comparison
for both architectures. The pattern is uniform: near-complete collapse in every character
group without exception.

\begin{table}[H]
\centering
\caption{Per-Group IoU Domain Gap (Attention U-Net)}
\label{tab:group_gap}
\begin{tabular}{lccc}
\toprule
\textbf{Group} & \textbf{Synthetic IoU} & \textbf{IAM IoU} & \textbf{Absolute Drop} \\
\midrule
Lowercase & 76.38\% & 0.37\% & $-76.01$ pp \\
Uppercase & 54.94\% & 0.04\% & $-54.90$ pp \\
Digits    & 49.51\% & 0.00\% & $-49.51$ pp \\
Special   & 47.74\% & 0.04\% & $-47.70$ pp \\
\bottomrule
\end{tabular}
\end{table}

\begin{table}[H]
\centering
\caption{Per-Group IoU Domain Gap (SwinUNet)}
\label{tab:group_gap_swin}
\begin{tabular}{lccc}
\toprule
\textbf{Group} & \textbf{Synthetic IoU} & \textbf{IAM IoU} & \textbf{Absolute Drop} \\
\midrule
Lowercase & 82.28\% & 0.70\% & $-81.58$ pp \\
Uppercase & 64.06\% & 0.00\% & $-64.06$ pp \\
Digits    & 56.69\% & 0.00\% & $-56.69$ pp \\
Special   & 51.78\% & 0.00\% & $-51.78$ pp \\
\bottomrule
\end{tabular}
\end{table}

Lowercase letters retain the highest residual IAM IoU in both models (0.37\% for U-Net,
0.70\% for SwinUNet). While negligible in absolute terms, this is consistent with two
reinforcing factors: lowercase characters are the most frequent class in English text and
therefore receive the most training examples, and the basic forms of common lowercase
letters (\eg `a', `e', `o', `n') are among the least idiosyncratic across writers, making
them somewhat more recognisable across the synthetic-to-real boundary. Digits and
uppercase letters in contrast collapse to exactly 0.00\% IAM IoU, indicating complete
failure of class-level discrimination for these groups on real handwriting.

\begin{figure}[htbp]
    \centering
    \begin{subfigure}[t]{0.48\textwidth}
        \centering
        \includegraphics[width=\textwidth]{figures/unet_iam_per_group_iou.png}
        \caption{U-Net --- per-group IoU}
        \label{fig:pg_iam_unet_iou}
    \end{subfigure}
    \hfill
    \begin{subfigure}[t]{0.48\textwidth}
        \centering
        \includegraphics[width=\textwidth]{figures/swin_iam_per_group_iou.png}
        \caption{SwinUNet --- per-group IoU}
        \label{fig:pg_iam_swin_iou}
    \end{subfigure}

    \vspace{0.6em}

    \begin{subfigure}[t]{0.48\textwidth}
        \centering
        \includegraphics[width=\textwidth]{figures/unet_iam_per_group_dice.png}
        \caption{U-Net --- per-group Dice}
        \label{fig:pg_iam_unet_dice}
    \end{subfigure}
    \hfill
    \begin{subfigure}[t]{0.48\textwidth}
        \centering
        \includegraphics[width=\textwidth]{figures/swin_iam_per_group_dice.png}
        \caption{SwinUNet --- per-group Dice}
        \label{fig:pg_iam_swin_dice}
    \end{subfigure}

    \vspace{0.6em}

    \begin{subfigure}[t]{0.48\textwidth}
        \centering
        \includegraphics[width=\textwidth]{figures/unet_iam_per_group_f1.png}
        \caption{U-Net --- per-group F1}
        \label{fig:pg_iam_unet_f1}
    \end{subfigure}
    \hfill
    \begin{subfigure}[t]{0.48\textwidth}
        \centering
        \includegraphics[width=\textwidth]{figures/swin_iam_per_group_f1.png}
        \caption{SwinUNet --- per-group F1}
        \label{fig:pg_iam_swin_f1}
    \end{subfigure}

    \caption{Per-group segmentation metrics on the IAM real handwriting test set for both
    architectures across IoU (top), Dice (middle), and F1 (bottom). Near-complete collapse
    is observed across every character group, with lowercase retaining a marginal edge
    consistent with its higher training frequency and lower writer-to-writer variability.
    The agreement of all three metrics confirms that the domain gap is not an artefact of
    any single measure.}
    \label{fig:per_group_iam}
\end{figure}

\begin{figure}[htbp]
    \centering
    \includegraphics[width=\textwidth]{figures/unet_iam_predictions.png}
    \caption{Attention U-Net predictions on IAM real handwriting samples (columns: input
    image, pseudo ground-truth mask, model prediction, colour overlay). The overlay
    confirms that foreground ink regions are detected in spatially correct locations,
    indicating that low-level ink-detection features transfer from synthetic training.
    However, the predicted class labels---visible as the colour pattern in the prediction
    column---do not match the pseudo ground truth. The model assigns ink pixels
    predominantly to high-frequency lowercase classes from the synthetic training
    distribution, irrespective of the actual character identity at each position. This
    behaviour explains the near-zero macro IoU alongside the reasonable pixel accuracy:
    the model has learned \emph{where} ink is but has not learned \emph{which character}
    that ink represents in the real handwriting domain.}
    \label{fig:iam_predictions}
\end{figure}

% =============================================================================
\section{Discussion}
\label{sec:discussion}
% =============================================================================

\subsection{Key Findings}

\subsubsection{Pixel-Level Masks Correct a Fundamental Training Objective Mismatch}

The shift from bounding-box fills to glyph-contour masks is the most structurally
significant change in this version, because it corrects what the model is trained to
predict. Version~1 assigned a rectangular fill to every character region: the model
optimised its predictions to match these boxes, which is a well-defined but wrong task.
The glyph-contour masks in version~2 assign class labels only to ink pixels, so the model
learns to predict the shape of the character as it actually appears.

The measurable consequence is visible in the validation IoU: 94.52\% in version~1 versus
98.08\% in version~2 for the Attention U-Net. It is not possible to attribute this gain to
the mask change alone, because the augmentation pipeline, class weighting scheme, and
present-class-only Dice averaging were all introduced simultaneously. However, the
expected contribution of each change can be reasoned about qualitatively. The mask quality
change is fundamental---it alters the ground truth that every gradient update optimises
toward, so its effect compounds across all training steps. The augmentation pipeline is a
regularisation improvement, primarily relevant to generalisation rather than peak
in-distribution accuracy. The class weighting and Dice averaging changes affect how
gradient signals are balanced across the 80 classes, improving the signal quality for
rare classes. On balance, the mask correction is expected to be the largest contributor
to the validation IoU improvement.

\subsubsection{SwinUNet Provides Better Class Discrimination on a Long-Tailed Vocabulary}

A counterintuitive result is that the model with lower validation IoU (SwinUNet, 95.10\%)
consistently outperforms the model with higher validation IoU (Attention U-Net, 98.08\%)
on every class-balanced test metric. This is not a contradiction---it is a direct
consequence of the metric inconsistency between training and evaluation described in
Section~\ref{sec:metrics}.

The deeper finding is that the two architectures optimise for qualitatively different
objectives under the same training protocol. The Attention U-Net's convolutional inductive
bias makes it excellent at learning precise local boundary features, which yields high
scores on pixel-weighted metrics dominated by the most frequent classes. SwinUNet's global
self-attention enables it to use long-range context---the surrounding characters and word
shape---when classifying each character position. This is exactly the kind of information
that disambiguates confusable character pairs (\eg `l' vs.\ `1') that appear at low
frequency and are penalised heavily in a class-balanced average.

This finding suggests that for long-tailed multi-class segmentation tasks, the choice of
evaluation metric is as consequential as the choice of architecture. A practitioner
optimising for pixel-weighted validation IoU would select the Attention U-Net; a
practitioner optimising for class-balanced macro IoU---which better reflects performance
across the full vocabulary---would select SwinUNet.

\subsubsection{The Synthetic-to-Real Gap is Severe and Structurally Two-Tiered}

The IAM results establish the most significant finding of this project: both architectures
learn two qualitatively distinct types of representation from synthetic training data, and
only one type transfers to real handwriting.

The first type is \emph{ink localisation}: discriminating ink pixels from background.
This is a low-level intensity-based task. The features that solve it---local intensity
gradients, edge responses, stroke boundaries---are largely invariant to the specific
character being written and to the difference between font-rendered and pen-drawn
appearance. Pixel accuracy degrades by only 12--15\% on real handwriting, confirming that
this type of representation transfers.

The second type is \emph{character identity}: assigning the correct class label to each
ink region. This requires recognising the specific shape signature of each character. The
features that solve it in the synthetic domain are trained on smooth, perfectly consistent
font-rendered glyphs, where the same character looks identical at every occurrence. Real
handwriting produces each character instance through a pen trajectory under variable
pressure, speed, and motor noise, resulting in local stroke variations, natural
irregularities at junctions and endpoints, and writer-specific idiosyncrasies that are
entirely absent from the training distribution. Character-level macro IoU collapses by
over 99\%, confirming that these identity features do not transfer at all.

The two-tier structure of the gap is itself a scientifically useful result. It precisely
isolates the bottleneck: the problem is not that the models fail to find the ink, but that
they cannot identify what character the ink represents once it is found. Any practical
solution must address this second layer specifically.

\subsection{Limitations}

\subsubsection{Approximate IAM Ground Truth}
\label{sec:approx_gt}

The pseudo ground-truth masks generated via connected component analysis are inherently
imperfect. Touching characters merge into a single component; multi-stroke characters
(\eg `i', `j', `t') produce multiple components that may not align with the transcription
sequence; and reading-order ambiguities arise at line boundaries. All of these introduce
noise into the evaluation masks. The reported IAM metrics consequently underestimate
actual model performance by an unknown amount. True IAM evaluation would require manual
pixel-level annotation, which does not currently exist for the database. The IAM numbers
should therefore be interpreted as lower bounds on real-data performance rather than exact
accuracy figures.

\subsubsection{Font Rendering vs.\ Natural Stroke Variation}

Even with the extended augmentation pipeline, the synthetic training data is fundamentally
font-rendered. Fonts define character shapes as smooth B\'{e}zier curves rendered
identically at every occurrence of the same character, in the same font, at the same size.
Real handwriting produces each character instance from a unique continuous pen trajectory,
yielding local stroke texture, natural irregularities at junctions and endpoints, and
ligatures between connected letters. The slant augmentation successfully introduces
cursive angle variation, and elastic deformation introduces global waviness, but neither
replicates the within-stroke microstructure that characterises handwritten glyphs. This is
the root cause of the domain gap and cannot be fully resolved through post-hoc image
augmentation of font-rendered output.

\subsubsection{Metric Inconsistency Between Training and Evaluation}
\label{sec:metric_inconsistency}

The model is optimised by early stopping for pixel-weighted validation IoU, but the
primary evaluation metric at test time is class-balanced macro IoU. These two quantities
can favour different model checkpoints. Concretely, the best checkpoint under
pixel-weighted validation IoU is not guaranteed to be the best checkpoint under
class-balanced macro IoU. This means the reported test results may not represent the best
achievable performance under the test metric---a model selected by a different criterion
could score higher on the class-balanced evaluation. Harmonising the training-time and
test-time metrics by periodically computing the full confusion-matrix macro IoU during
training, and using it directly as the early stopping criterion, would address this
inconsistency and is a straightforward improvement requiring no change to architecture or
data.

\subsection{Future Work}

\subsubsection{Domain Adaptation}

The most impactful next step is to directly address the synthetic-to-real gap at the
representation level. Adversarial domain adaptation---training a domain discriminator
alongside the segmentation model to align feature distributions between synthetic and real
images---could transfer character-identity knowledge learned from synthetic data to real
handwriting without requiring real pixel-level annotations. Alternatively, collecting even
a modest number of manually annotated real handwriting lines (100--200 images) for
fine-tuning could close much of the gap through targeted supervised learning, leveraging
the strong synthetic pre-training as initialisation.

\subsubsection{Stroke-Level Synthetic Generation}

Replacing font rendering with a stroke-level handwriting synthesis model would produce
training data that more closely mirrors the real handwriting distribution at the source.
Models that generate pen trajectories conditioned on character identity produce diverse,
naturalistic character instances whose within-stroke microstructure is far closer to what
a human writer produces. This would reduce the domain gap at the data level rather than
attempting to close it through augmentation or adaptation after the fact.

\subsubsection{Hybrid Architecture}

The complementary strengths observed in this comparison---U-Net's precise local boundary
detection versus SwinUNet's global class discrimination---motivate a hybrid design
combining a CNN encoder with a transformer decoder. Architectures such as TransUNet or
HiFormer explore exactly this combination. In the context of this task, a CNN encoder
could provide rich local stroke features to the decoder, while a transformer decoder could
use global context across the character sequence to resolve class ambiguity.

\subsubsection{Harmonised Evaluation}

Aligning the training-time validation metric with the test-time evaluation metric would
ensure that model selection via early stopping optimises for the quantity reported at
evaluation. Running the full confusion-matrix class-balanced macro IoU periodically during
training---rather than only at test time---would provide a more reliable basis for both
early stopping and hyperparameter tuning, and would likely further improve SwinUNet's
already superior class-balanced test performance.

% =============================================================================
\section{Conclusion}
\label{sec:conclusion}
% =============================================================================

This project directly addressed all three improvements requested following the version~1
review---pixel-level glyph-contour masks, a comprehensive augmentation pipeline with
explicit real-handwriting motivation for each transform, and quantitative evaluation on
real handwriting from the IAM database---and extended the work with a controlled
two-architecture comparison between Attention U-Net and SwinUNet.

Three findings stand out. First, the glyph-contour masks correct a fundamental training
objective mismatch from version~1, contributing to a validation IoU improvement from
94.52\% to 98.08\% for the Attention U-Net. Second, the architecture comparison reveals
a meaningful trade-off: Attention U-Net achieves higher pixel-weighted validation IoU,
while SwinUNet achieves higher class-balanced macro IoU on all test metrics, suggesting
that global self-attention better handles the long-tailed 80-class vocabulary. Third, and
most significantly, the IAM evaluation reveals a structurally two-tiered domain gap:
pixel accuracy holds at 84--87\% on real handwriting, confirming partial transfer of
low-level ink-detection features, while character-class macro IoU collapses below 0.3\%,
confirming near-total failure of character-identity features to generalise beyond the
synthetic domain.

This last finding is a precisely characterised negative result with direct scientific
value: it establishes a clear, reproducible lower bound on the current approach and
isolates exactly where the failure occurs. The models are effective text detectors but not
character classifiers when evaluated on real handwriting. Addressing this gap requires
either closing the distance between synthetic and real training distributions (through
stroke-level synthesis or domain adaptation) or providing a modest amount of real
annotated data for fine-tuning. The complete codebase, with centralised configuration,
W\&B experiment tracking, and separate evaluation tools for synthetic and real data, is
structured to support exactly these next steps.

% =============================================================================
% References
% =============================================================================

\newpage
\begin{thebibliography}{10}

\bibitem{ronneberger2015}
O.~Ronneberger, P.~Fischer, and T.~Brox,
``U-Net: Convolutional Networks for Biomedical Image Segmentation,''
\textit{Medical Image Computing and Computer-Assisted Intervention (MICCAI)},
vol.~9351, pp.~234--241, 2015.

\bibitem{oktay2018}
O.~Oktay, J.~Schlemper, L.~L.~Folgoc, \emph{et al.},
``Attention U-Net: Learning Where to Look for the Pancreas,''
\textit{Medical Imaging with Deep Learning (MIDL)}, 2018.

\bibitem{cao2022}
H.~Cao, Y.~Wang, J.~Chen, \emph{et al.},
``Swin-Unet: Unet-like Pure Transformer for Medical Image Segmentation,''
\textit{ECCV Workshop on Computer Vision for Automated Medical Diagnosis}, 2022.

\bibitem{liu2021}
Z.~Liu, Y.~Lin, Y.~Cao, \emph{et al.},
``Swin Transformer: Hierarchical Vision Transformer using Shifted Windows,''
\textit{IEEE/CVF International Conference on Computer Vision (ICCV)},
pp.~10012--10022, 2021.

\bibitem{marti2002}
U.-V.~Marti and H.~Bunke,
``The IAM-database: An English Sentence Database for Offline Handwriting Recognition,''
\textit{International Journal on Document Analysis and Recognition},
vol.~5, no.~1, pp.~39--46, 2002.

\bibitem{loshchilov2019}
I.~Loshchilov and F.~Hutter,
``Decoupled Weight Decay Regularization,''
\textit{International Conference on Learning Representations (ICLR)}, 2019.

\bibitem{sudre2017}
C.~H.~Sudre, W.~Li, T.~Vercauteren, S.~Ourselin, and M.~J.~Cardoso,
``Generalised Dice Overlap as a Deep Learning Loss Function for Highly Unbalanced
Segmentations,''
\textit{Deep Learning in Medical Image Analysis (DLMIA)},
vol.~10553, pp.~240--248, 2017.

\bibitem{paszke2019}
A.~Paszke, S.~Gross, F.~Massa, \emph{et al.},
``PyTorch: An Imperative Style, High-Performance Deep Learning Library,''
\textit{Advances in Neural Information Processing Systems (NeurIPS)},
vol.~32, pp.~8024--8035, 2019.

\bibitem{falcon2019}
W.~Falcon and {The PyTorch Lightning team},
``PyTorch Lightning,''
version~2.1.4, available at \url{https://github.com/Lightning-AI/pytorch-lightning},
2019.

\bibitem{wandb2020}
L.~Biewald,
``Experiment Tracking with Weights and Biases,''
Software available from \url{https://wandb.com}, 2020.

\bibitem{albumentations2020}
A.~Buslaev, V.~I.~Iglovikov, E.~Khvedchenya, \emph{et al.},
``Albumentations: Fast and Flexible Image Augmentations,''
\textit{Information}, vol.~11, no.~2, p.~125, 2020.

\end{thebibliography}

% =============================================================================
% Declaration of Authorship
% =============================================================================

\newpage
\thispagestyle{empty}
\section*{Declaration of Authorship}

I confirm that I have written this report unaided and without using sources other than
those listed, and that this report has never been submitted to another examination
authority and accepted as part of an examination achievement, neither in this form nor in
a similar form. All content taken from a third party, either verbatim or in substance,
has been acknowledged as such.

\vspace{3cm}

\noindent
Erlangen, 18th~March~2026 \hfill Fahad Tariq

\end{document}